\documentclass[conference]{IEEEtran}

\usepackage[utf8]{inputenc}
\usepackage[T1]{fontenc}
\usepackage[brazil]{babel}
\usepackage{amsmath}
\usepackage{graphicx}
\usepackage{url}
\usepackage{hyperref}
\usepackage{listings}
\usepackage{xcolor}
\usepackage{booktabs}
\usepackage{array}
\usepackage{multirow}
\usepackage{float}
\usepackage{cite}

\definecolor{codebg}{HTML}{F6F8FA}
\definecolor{codeframe}{HTML}{D0D7DE}
\definecolor{codetext}{HTML}{24292E}
\definecolor{codecomment}{HTML}{6E7781}
\definecolor{codekeyword}{HTML}{CF222E}
\definecolor{codestring}{HTML}{0A3069}
\definecolor{codenumber}{HTML}{0550AE}

\lstset{
  backgroundcolor=\color{codebg},
  basicstyle=\ttfamily\scriptsize\color{codetext},
  commentstyle=\color{codecomment}\itshape,
  keywordstyle=\color{codekeyword}\bfseries,
  stringstyle=\color{codestring},
  numberstyle=\tiny\color{codenumber},
  breaklines=true,
  frame=single,
  rulecolor=\color{codeframe},
  numbers=left,
  numbersep=6pt,
  tabsize=4,
  showspaces=false,
  showstringspaces=false,
  xleftmargin=12pt,
  xrightmargin=4pt,
  language=Python
}

\hypersetup{
  colorlinks=true,
  linkcolor=blue,
  urlcolor=blue,
  citecolor=blue
}

\begin{document}

\title{CriptoIA: Avaliação Empírica de Recomendações Multimodais em Criptomoedas\\com Backtesting Sistemático, Comparação com LSTM e Interface Web Segura}

\author{
  \IEEEauthorblockN{Gabriel Teixeira}
  \IEEEauthorblockA{
    Ciência da Computação\\
    Universidade Franciscana\\
    Santa Maria, Brasil\\
    gabriel.teixeira@ufn.edu.br
  }
  \and
  \IEEEauthorblockN{Orientador: Prof. Mirkos O. Martins}
  \IEEEauthorblockA{
    Ciência da Computação\\
    Universidade Franciscana\\
    Santa Maria, Brasil\\
    mirkos@ufn.edu.br
  }
}

\maketitle

\begin{abstract}
Este trabalho apresenta a etapa final do CriptoIA, um sistema de análise de criptomoedas que integra coleta automatizada de dados da Binance, cálculo de indicadores técnicos em Pandas, análise multimodal com o modelo Google Gemini e backtesting sistemático. No TFG II, foram implementados três experimentos complementares de avaliação objetiva: (i) backtesting de uma estratégia proxy determinística baseada em MM20, MM50, RSI-14, MACD e Bandas de Bollinger, cobrindo 4 ativos (BTC, ETH, SOL, DOGE) × 4 timeframes (1h, 4h, 1d, 1w) ao longo de 18 meses (jan/2025–jun/2026); (ii) backtesting multimodal real do Gemini 3.5 Flash-Lite no timeframe diário, com amostragem semanal (63 pontos por ativo, 252 inferências totais, janela de 90 candles passados, zero fallbacks); e (iii) complemento exploratório intraday do Gemini em 1h e 4h (jul–out/2026, 304 inferências, zero fallbacks). Também foi desenvolvida comparação com Buy-and-Hold e LSTM real em PyTorch (walk-forward, mesma janela out-of-sample), interface web Flask com autenticação, CSRF, histórico persistido e visualização dos resultados, e suíte de 200 testes automatizados com 92\% de cobertura. Os critérios de sucesso do TCC (acurácia direcional > 50\% em pelo menos 3 de 4 ativos e retorno não negativo) não foram atendidos em nenhum dos três recortes: no Experimento 1 (proxy, 1d), 2/4 ativos atenderam acurácia e 1/4 retorno; no Experimento 2 (Gemini real, 1d), 3/4 atenderam acurácia mas 0/4 retorno; no Experimento 3 (Gemini 1h/4h), 7/8 combinações atenderam acurácia mas apenas 3/8 retorno não negativo e 1/8 superou Buy-and-Hold. A contribuição está na metodologia reproduzível, na transparência dos resultados negativos e na integração completa de pipeline, backtesting, comparação, testes e interface — sem números fabricados.
\end{abstract}

\begin{IEEEkeywords}
Criptomoedas, Análise Técnica, Inteligência Artificial Multimodal, Google Gemini, Backtesting, LSTM, PyTorch, Walk-Forward, Python, CCXT, Plotly, Flask.
\end{IEEEkeywords}

\section{Introdução}

O mercado de criptomoedas consolidou-se como segmento financeiro relevante: em 2024, a capitalização total ultrapassou dois trilhões de dólares e, no Brasil, a Receita Federal registrou volumes recordes de declarações envolvendo ativos digitais \cite{coinmarketcap2024,pinheiro2019mercado}. As bolsas operam 24/7 sem interrupção, com volatilidade intrínseca que torna o acompanhamento manual contínuo inviável para a maioria dos investidores \cite{assaf2012mercado}. Nesse cenário, ferramentas de análise técnica automatizada surgem como recurso valioso para identificar padrões e oportunidades de forma sistemática.

A análise técnica fundamenta-se no estudo de dados históricos de preço e volume \cite{murphy1999technical}. Indicadores como Médias Móveis Simples (SMA), RSI (\textit{Relative Strength Index}), MACD (\textit{Moving Average Convergence Divergence}) e Bandas de Bollinger são amplamente utilizados para identificar tendências, \textit{momentum}, sobrecompra e sobrevenda \cite{cavalcante2005mercado,wilder1978new,bollinger2001bollinger}. Paralelamente, modelos de linguagem grandes com capacidade multimodal (\textit{Vision-Language Models}) abrem novas possibilidades: o Gemini \cite{google2024gemini} processa simultaneamente texto e imagens, permitindo interpretar gráficos financeiros de forma visual e incorporar padrões gráficos difíceis de capturar apenas por métricas numéricas.

No TFG I, foi proposto e desenvolvido o \textbf{CriptoIA} — sistema que integra coleta de dados da Binance via CCXT, cálculo de cinco indicadores técnicos em Pandas, geração de gráfico técnico PNG e análise multimodal pelo Gemini 2.5 Flash, com saída estruturada em JSON e dashboard interativo Plotly \cite{tcc1}. A hipótese central era que o enriquecimento do \textit{prompt} com interpretações qualitativas automáticas dos indicadores, combinado ao envio simultâneo do gráfico em alta resolução, permitiria recomendações mais fundamentadas do que uma consulta genérica sem contexto especializado.

O TFG II tem como objetivo principal a \textbf{validação objetiva} dessa hipótese por meio de backtesting sistemático, comparação com baselines estabelecidas e entrega de uma interface web segura. As questões de pesquisa que guiam esta etapa são:

\begin{enumerate}
  \item \textit{Qual a acurácia direcional e o retorno simulado das recomendações geradas pelo CriptoIA (proxy determinístico) e pelo Gemini multimodal real, em comparação com Buy-and-Hold e LSTM, ao longo de 18 meses e 4 timeframes?}
  \item \textit{A integração multimodal (gráfico + indicadores) produz recomendações superiormente fundamentadas em relação a estratégias puramente quantitativas ou puramente visuais?}
  \item \textit{É possível entregar uma interface web reproduzível, testável e segura para execução e visualização desses experimentos?}
\end{enumerate}

As contribuições específicas deste trabalho (TFG II) são:

\begin{itemize}
  \item \textbf{Três experimentos de backtesting complementares e reproduzíveis:} (1) proxy determinístico 4×4 (16 combinações, 18 meses); (2) Gemini multimodal real em 1d amostral (252 inferências, zero fallbacks); (3) complemento exploratório Gemini em 1h/4h (304 inferências, zero fallbacks);
  \item \textbf{Comparação sistemática com baselines:} Buy-and-Hold passivo e LSTM real em PyTorch com walk-forward e janela comum out-of-sample;
  \item \textbf{Interface web Flask completa:} autenticação com hash, CSRF, rate limit, sessões seguras, histórico persistido atomicamente e visualização dinâmica dos três experimentos sem números fixos;
  \item \textbf{Suite de testes abrangente:} 200 testes (CT01–CT06 + módulos) com 92\% de cobertura, mocks de rede, sem dependência de internet;
  \item \textbf{Metodologia anti \textit{look-ahead bias}:} janela de 90 candles estritamente passados, sinal em $t$ aplicado a retorno em $t+1$, cache atômico de inferências Gemini para retomada sem duplicação.
\end{itemize}

Este artigo está organizado da seguinte forma: a Seção~\ref{sec:correlatos} atualiza os trabalhos correlatos; a Seção~\ref{sec:referencial} complementa o referencial teórico com conceitos de validação experimental; a Seção~\ref{sec:metodologia} detalha a metodologia implementada no TFG II; a Seção~\ref{sec:implementacao} descreve os módulos implementados; a Seção~\ref{sec:resultados} apresenta e discute os resultados dos três experimentos; e a Seção~\ref{sec:conclusao} conclui com limitações e trabalhos futuros.

\section{Trabalhos Correlatos Atualizados}
\label{sec:correlatos}

Além dos três trabalhos brasileiros analisados no TFG I \cite{monteiro2026ssinnova,silvajunior2022allbot,franciscon2023lstm}, dois estudos recentes internacionais contextualizam a validação empírica de sistemas baseados em LLM para finanças:

\textbf{Liu et al. (2024) \cite{zhang2024finben}} propuseram o \textit{FinBen}, framework de avaliação de LLMs em tarefas financeiras (análise de sentimento, previsão de movimento de preço, geração de relatórios). Demonstraram que modelos genéricos sem \textit{fine-tuning} especializado apresentam desempenho inferior a modelos dedicados, reforçando a importância de \textit{prompt engineering} estruturado e contexto técnico — princípios centrais do CriptoIA.

\textbf{Yang et al. (2024) \cite{yang2024fingpt}} desenvolveram o \textit{FinGPT}, modelo financeiro open-source, e \textbf{Lopez-Lira e Tang (2024) \cite{lopez2024chatgpt}} realizaram backtesting de sinais gerados por ChatGPT em ações dos EUA, comparando com estratégias técnicas clássicas. Encontraram que a acurácia direcional do LLM superou heurísticas simples em cenários de alta volatilidade, mas o retorno líquido (após custos) não diferiu significativamente de Buy-and-Hold. O estudo destaca a necessidade de relatar retornos compostos com sobreposição de horizontes e custos de transação — limitações explicitamente declaradas no CriptoIA (Seção~\ref{sec:conclusao}).

Esses trabalhos confirmam que a abordagem do CriptoIA — \textit{prompt engineering} com interpretações qualitativas, janela de contexto fixa, validação via backtesting com baselines múltiplas e relato transparente de limitações — está alinhada com o estado da arte em avaliação empírica de LLMs para finanças.

\section{Referencial Teórico Complementar}
\label{sec:referencial}

\subsection{Validação Experimental em Finanças Computacionais}

A avaliação de sistemas de recomendação financeira exige metodologia rigorosa para evitar vieses comuns \cite{bailey2014pseudomathematics}:

\begin{itemize}
  \item \textbf{\textit{Look-ahead bias}:} uso de informação futura na geração do sinal. O CriptoIA garante que, em cada ponto de decisão $t$, apenas os 90 candles até $t$ (inclusive) são usados para calcular indicadores e gerar o gráfico enviado ao modelo.
  \item \textbf{\textit{Survivorship bias}:} avaliação apenas em ativos que sobreviveram. O CriptoIA fixa os 4 pares \textit{a priori} (BTC, ETH, SOL, DOGE) antes da coleta.
  \item \textbf{\textit{Data snooping}:} otimização implícita de hiperparâmetros no mesmo conjunto de teste. O proxy CriptoIA usa regras fixas (MM20 > MM50, RSI < 70, MACD > Sinal) sem otimização; o LSTM usa \textit{walk-forward} com re-treino periódico apenas com dados até $t$.
  \item \textbf{\textit{Overlapping horizons}:} horizontes de avaliação que se sobrepõem, inflando a contagem de observações independentes. Nos Experimentos 2 e 3, o passo de amostragem (7 dias em 1d, 3 dias em 1h/4h) pode gerar sobreposição com horizontes de 7–14 dias; o retorno é reportado como \textit{composto por decisão}, não como curva de capital sem sobreposição (limitação declarada).
\end{itemize}

\subsection{LSTM com Walk-Forward}

Redes LSTM (\textit{Long Short-Term Memory}) são arquiteturas recorrentes capazes de capturar dependências temporais de longo prazo \cite{hochreiter1997lstm}. Em finanças, o \textit{walk-forward} (também chamado \textit{rolling window} ou \textit{expanding window}) é o padrão para avaliação temporal honesta \cite{fischer2018deep}: o modelo é treinado em uma janela inicial, prevê o próximo passo, a janela avança, o modelo é re-treinado (ou atualizado) e o processo repete. O CriptoIA implementa \textit{walk-forward} com re-treino a cada \texttt{passo\_retreino} candles, usando apenas dados até o instante de decisão $t$ para treino e validação, eliminando vazamento temporal.

\subsection{Interface Web Segura para Sistemas Acadêmicos}

Aplicações web expostas a credenciais e dados sensíveis devem implementar: autenticação com \textit{hash} de senha (PBKDF2/SHA-256 via Werkzeug), proteção CSRF (\textit{Cross-Site Request Forgery}) via token síncrono \cite{owasp2021csrf}, cookies \texttt{HttpOnly}, \texttt{SameSite=Lax} e opção \texttt{Secure} para HTTPS, rate limiting contra força bruta, e headers de segurança (\texttt{X-Content-Type-Options}, \texttt{X-Frame-Options}, \texttt{Referrer-Policy}) \cite{owasp2021headers}. O CriptoIA implementa todos esses controles no módulo \texttt{webapp/}.

\section{Metodologia do TFG II}
\label{sec:metodologia}

\subsection{Período, Ativos e Timeframes}

A Tabela~\ref{tab:periodo} resume o escopo experimental. O período principal (Experimentos 1 e 2) cobre 18 meses (jan/2025–jun/2026), capturando ciclos de alta e baixa. O Experimento 3 utiliza janela exploratória de 4 meses (jul–out/2026) para complementar a evidência intraday sem alterar o recorte oficial.

\begin{table}[H]
  \centering
  \caption{Escopo experimental do TFG II}
  \label{tab:periodo}
  \small
  \begin{tabular}{lll}
    \toprule
    \textbf{Dimensão} & \textbf{Experimentos 1 e 2} & \textbf{Experimento 3} \\
    \midrule
    Período         & 2025-01-01 a 2026-06-30   & 2026-07-01 a 2026-10-31 \\
    Pares           & BTC, ETH, SOL, DOGE/USDT  & mesmos 4 pares          \\
    Timeframes      & 1h, 4h, 1d, 1w (Exp. 1)   & 1h, 4h                  \\
                    & 1d apenas (Exp. 2)        &                         \\
    Modelo Gemini   & \texttt{gemini-3.5-flash-lite} (configurável) & mesmo modelo \\
    \bottomrule
  \end{tabular}
\end{table}

A API pública da Binance limita a 1.000 candles por requisição; a coleta implementa paginação automática, deduplicação por \texttt{timestamp} e marca resultados incompletos com atributos de auditoria (\texttt{coleta\_incompleta}, \texttt{erro\_coleta}).

\subsection{Estratégia Proxy CriptoIA (Experimento 1)}

A estratégia determinística proxy replica a lógica técnica que fundamenta o \textit{prompt} enviado ao Gemini, permitindo avaliação em grade 4×4 completa sem consumo de API:

\begin{align}
\text{COMPRAR} &: \text{MM20} > \text{MM50} \land \text{RSI}_{14} < 70 \land \text{MACD} > \text{Sinal\_MACD} \\
\text{VENDER}  &: \text{MM20} < \text{MM50} \land \text{RSI}_{14} > 30 \land \text{MACD} < \text{Sinal\_MACD} \\
\text{MANTER}   &: \text{demais casos}
\end{align}

É uma estratégia \textit{long-only}: sinal $+1$ (COMPRAR) mantém posição comprada; $0$ ou $-1$ (VENDER/MANTER) mantém em caixa. O sinal em $t$ é aplicado ao retorno futuro $r_{t+1} = \frac{\text{close}_{t+1}}{\text{close}_t} - 1$ (\textit{shift}(-1)), nunca usado para gerar o sinal. A acurácia direcional é calculada apenas sobre sinais $\pm 1$ (exclui MANTER).

\subsection{Backtesting Multimodal Real do Gemini (Experimentos 2 e 3)}

A Figura~\ref{fig:gemini_backtest} ilustra o protocolo. Em cada ponto de decisão $t$ (amostrado sistematicamente):

\begin{enumerate}
  \item Gráfico PNG gerado com \textbf{apenas} os 90 candles até $t$ (sem dados futuros na imagem nem nos indicadores).
  \item Indicadores calculados sobre essa mesma janela de 90 candles.
  \item \textit{Prompt} estruturado + imagem enviados ao Gemini via API multimodal.
  \item Resposta JSON validada: \texttt{acao} $\in$ \{COMPRAR, VENDER, MANTER\}, \texttt{justificativa}, \texttt{preco\_alvo}, \texttt{tempo\_estimado}.
  \item Em falha persistente (até 3 tentativas): fallback \texttt{MANTER} com \texttt{fallback=True}.
  \item \texttt{tempo\_estimado} convertido em janela de dias (ex.: ``7 a 14 dias'' $\rightarrow$ 7–14; ``2 a 4 semanas'' $\rightarrow$ 14–28).
  \item \texttt{preco\_alvo} normalizado (formatos BR/US e percentual) e comparado com máximas/mínimas futuras na janela.
  \item Cache atômico em \texttt{resultados/gemini\_backtest\_cache.json} permite retomada sem chamadas duplicadas.
\end{enumerate}

\begin{figure}[H]
  \centering
  \includegraphics[width=\columnwidth]{imagens/diagrama_fluxo.png}
  \caption{Fluxo do backtesting multimodal do Gemini: janela de 90 candles passados, inferência no ponto $t$, avaliação no horizonte estimado.}
  \label{fig:gemini_backtest}
\end{figure}

Métricas calculadas por par/timeframe: acurácia em $t+1$, acurácia no prazo estimado, taxa de atingimento do alvo, retorno composto por decisão, Buy-and-Hold na mesma janela e concordância ponto a ponto com o proxy.

\subsection{Comparação com Buy-and-Hold e LSTM (Walk-Forward)}

O LSTM em PyTorch utiliza 5 \textit{features}: RSI normalizado, diferença MACD/Sinal, distância relativa MM20/MM50, posição nas Bandas de Bollinger e retorno do candle anterior. Alvo: direção do próximo candle (binário). \textit{StandardScaler} ajustado apenas no treino. Arquitetura: LSTM (1 camada, \texttt{hidden\_size}=32) + Dropout (0.2) + Linear. Treino: 20 épocas, \texttt{seq\_len}=24, \textit{seed}=42, CPU only. Walk-forward com re-treino a cada 1.000 candles. No timeframe 1w, \texttt{seq\_len} reduzido automaticamente para 7 e registrado.

A \textbf{janela comum} é o período out-of-sample do LSTM (últimos 40\% dos dados após treino inicial). As três estratégias (CriptoIA proxy, Buy-and-Hold, LSTM) são recalculadas exatamente nessa janela, eliminando comparação injusta entre períodos diferentes.

\subsection{Critérios de Sucesso do TCC}

\begin{itemize}
  \item \textbf{Critério 1 (Acurácia):} acurácia direcional de prazo > 50\% em pelo menos 3 dos 4 ativos.
  \item \textbf{Critério 2 (Retorno):} retorno acumulado simulado não negativo no período de avaliação.
  \item \textbf{Conjunto:} ambos os critérios devem ser atendidos simultaneamente.
\end{itemize}

O TCC I define os critérios por ativo, mas não especifica o timeframe do resumo. Para não criar parâmetro inexistente, preservam-se os 16 resultados brutos da grade 4×4 e adota-se o timeframe diário apenas como \textbf{recorte operacional} do resumo (quando disponível), explicitando essa escolha em \texttt{nota\_timeframe} nos artefatos JSON.

\subsection{Interface Web e Segurança}

A aplicação Flask (\texttt{webapp/}) entrega:
\begin{itemize}
  \item Login com credenciais por variáveis de ambiente (\texttt{WEB\_ADMIN\_USER}, \texttt{WEB\_ADMIN\_PASS}), senha com \textit{hash} PBKDF2:SHA-256, senha padrão apenas se variável não definida.
  \item CSRF em todos POSTs (token síncrono \texttt{hmac.compare\_digest}).
  \item Cookies \texttt{HttpOnly}, \texttt{SameSite=Lax}, \texttt{Secure} opcional.
  \item Rate limit: 5 tentativas/IP/5 min (em memória; produção requer Redis).
  \item Sessão renovada no login (\texttt{session.clear()} antes de autenticar).
  \item Rotas protegidas por \texttt{@login\_required}.
  \item Validação de par (\texttt{regex}) e timeframe; limite de quantidade 50–1000.
  \item Histórico de análises e backtests persistido atomicamente (\texttt{history\_store.py}: \texttt{fcntl}/\texttt{msvcrt} lock, escrita atômica via arquivo temporário + \texttt{os.replace}).
  \item Página de backtesting carrega JSONs reais dos 3 experimentos; resumo dinâmico dos critérios do TCC; sem tabelas acadêmicas fixas.
\end{itemize}

\subsection{Testes Automatizados}

Suite \texttt{pytest} com 200 testes:
\begin{itemize}
  \item CT01–CT06: requisitos funcionais e não funcionais do TCC I.
  \item Testes de módulos: backtesting, backtesting Gemini, comparação, múltiplos ativos, histórico, autenticação.
  \item Testes de integração: pipeline completo com mocks de rede e diretórios temporários.
  \item Cobertura: 92\% total (1.660 statements, 126 não cobertos); \texttt{analise\_tecnica.py}, \texttt{gerador\_grafico.py}, \texttt{multiplos\_ativos.py}, \texttt{gerador\_dashboard.py} com 100\%.
  \item Zero dependência de internet ou APIs reais; mocks via \texttt{unittest.mock} e \texttt{monkeypatch}.
\end{itemize}

\section{Implementação dos Módulos do TFG II}
\label{sec:implementacao}

\subsection{Novos Módulos do Backend}

\subsubsection{\texttt{backtesting.py} — Experimento 1 (Proxy Determinístico)}

Implementa coleta paginada para períodos longos, cálculo dos sinais proxy, métricas anti \textit{look-ahead}, execução paralela (\texttt{ThreadPoolExecutor}) para a grade 4×4, critérios de sucesso do TCC, exportação JSON/CSV e relatório textual. A função \texttt{executar\_backtesting\_completo()} orquestra a grade e retorna lista de dicionários com métricas por combinação par × timeframe.

\subsubsection{\texttt{backtesting\_gemini.py} — Experimentos 2 e 3 (Gemini Real)}

Módulo mais complexo do TFG II. Implementa:
\begin{itemize}
  \item Amostragem configurável (\texttt{passo\_amostragem} em dias, convertido para candles por timeframe).
  \item Janela fixa de 90 candles passados para gráfico e indicadores.
  \item Parser robusto de \texttt{tempo\_estimado} (intervalos, semanas, meses, texto livre) e \texttt{preco\_alvo} (BR/US, percentual, absoluto).
  \item Cache atômico com \textit{file locking} e retomada automática.
  \item Buy-and-Hold calculado na mesma janela de cada decisão.
  \item Concordância ponto a ponto com o proxy (mesma janela).
  \item Duas funções de entrada: \texttt{executar\_backtest\_gemini\_par()} (único par) e \texttt{executar\_backtest\_gemini\_grade()} (grade completa com \texttt{--intervalos}).
\end{itemize}

\subsubsection{\texttt{comparacao.py} — Buy-and-Hold e LSTM}

Implementa Buy-and-Hold passivo, extração de \textit{features}/alvo, LSTM PyTorch real, \textit{walk-forward} com \texttt{StandardScaler} apenas no treino, janela comum out-of-sample, exportação JSON/CSV e relatório comparativo. A classe \texttt{LSTMPredictor} encapsula modelo, escalonador e loop de treino/predição.

\subsubsection{\texttt{multiplos\_ativos.py} — Análise Paralela}

Usa \texttt{ThreadPoolExecutor} para analisar $N$ pares × $M$ timeframes em paralelo, com isolamento de erros (falha em um não interrompe os demais). Gera relatório consolidado em texto, JSON e CSV (\texttt{salvar\_relatorio\_csv}).

\subsection{Interface Web (\texttt{webapp/})}

\begin{itemize}
  \item \texttt{app.py}: rotas \texttt{/login}, \texttt{/logout}, \texttt{/} (dashboard), \texttt{/analise}, \texttt{/historico}, \texttt{/backtesting}. Carrega relatórios JSON consolidados dos 3 experimentos e renderiza templates Jinja2.
  \item \texttt{auth.py}: \texttt{verificar\_credenciais()}, \texttt{login\_required} decorator, \texttt{USANDO\_SENHA\_PADRAO} flag.
  \item \texttt{history\_store.py}: persistência JSON atômica com \textit{lock}, limites de 100 análises e 50 backtests.
  \item Templates: \texttt{login.html}, \texttt{dashboard.html}, \texttt{analise.html}, \texttt{historico.html}, \texttt{backtesting.html} (com abas para Experimento 1, 2 e 3).
  \item Estáticos: \texttt{style.css} (tema escuro, responsivo, animações).
\end{itemize}

\subsection{Orquestrador (\texttt{principal.py})}

Novos modos CLI:
\begin{itemize}
  \item \texttt{--backtesting} [\texttt{--grade}]: Experimento 1.
  \item \texttt{--backtesting-gemini} [\texttt{--grade}] [\texttt{--intervalos}] [\texttt{--passo}] [\texttt{--delay}] [\texttt{--prefixo-saida}]: Experimentos 2 e 3.
  \item \texttt{--comparacao} [\texttt{--grade}]: comparação com LSTM e Buy-and-Hold.
  \item \texttt{--multiplos} [\texttt{--pares}] [\texttt{--timeframes}] [\texttt{--usar-ia}] [\texttt{--salvar-csv}]: análise paralela.
  \item Parâmetros de data (\texttt{--inicio}, \texttt{--fim}), LSTM (\texttt{--epochs}, \texttt{--seq-len}, \texttt{--hidden-size}, \texttt{--passo-retreino}) e workers.
\end{itemize}

Fallback estruturado \texttt{MANTER} em qualquer exceção na pipeline, persistido em \texttt{resultados/fallback\_analise.json}.

\section{Resultados e Discussão}
\label{sec:resultados}

\subsection{Experimento 1 — Proxy Determinístico (Grade 4×4)}

A Tabela~\ref{tab:exp1_grade} apresenta os resultados da estratégia proxy nos 4 ativos × 4 timeframes. A cobertura foi de 100\% do período estimado nas 16 combinações.

\begin{table*}[t]
  \centering
  \caption{Experimento 1 — Proxy CriptoIA (jan/2025–jun/2026)}
  \label{tab:exp1_grade}
  \small
  \begin{tabular}{lcccc}
    \toprule
    Par & TF & Acurácia & Retorno & Max. DD \\
    \midrule
    BTC/USDT & 1h & 48.5\% & -1.67\%  & -15.30\% \\
    BTC/USDT & 4h & 49.0\% & -4.72\%  & -16.54\% \\
    BTC/USDT & 1d & 56.5\% & \textbf{+16.20\%} & -14.08\% \\
    BTC/USDT & 1w & 33.3\% & 0.00\%   & 0.00\%   \\
    ETH/USDT & 1h & 47.6\% & -23.26\% & -26.67\% \\
    ETH/USDT & 4h & 47.9\% & +5.96\%  & -17.69\% \\
    ETH/USDT & 1d & 47.4\% & -15.83\% & -28.70\% \\
    ETH/USDT & 1w & 50.0\% & 0.00\%   & 0.00\%   \\
    SOL/USDT & 1h & 48.9\% & -11.79\% & -38.02\% \\
    SOL/USDT & 4h & 47.6\% & -1.10\%  & -28.14\% \\
    SOL/USDT & 1d & 49.6\% & -39.06\% & -39.85\% \\
    SOL/USDT & 1w & N/A    & 0.00\%   & 0.00\%   \\
    DOGE/USDT& 1h & 48.8\% & +1.78\%  & -34.10\% \\
    DOGE/USDT& 4h & 46.7\% & +11.55\% & -30.48\% \\
    DOGE/USDT& 1d & 55.2\% & -5.08\%  & -41.04\% \\
    DOGE/USDT& 1w & 33.3\% & 0.00\%   & 0.00\%   \\
    \bottomrule
  \end{tabular}
\end{table*}

No recorte diário (recorte operacional do resumo):

\begin{itemize}
  \item Acurácia > 50\%: \textbf{2 de 4} ativos (BTC 56.5\%, DOGE 55.2\%; ETH 47.4\%, SOL 49.6\%).
  \item Retorno não negativo: \textbf{1 de 4} ativos (BTC +16.20\%; ETH -15.83\%, SOL -39.06\%, DOGE -5.08\%).
  \item \textbf{Critérios do TCC não atendidos.}
\end{itemize}

Observações: em 1w a estratégia \textit{long-only} não encontrou sinais de compra (apenas venda, que não abre posição vendida), resultando em retorno 0\%. A amostra semanal tem apenas 28 candles válidos após \textit{warm-up}.

\subsection{Experimento 2 — Gemini Multimodal Real (1d, Amostral)}

Modelo \texttt{gemini-3.5-flash-lite}; 63 pontos por par (252 inferências totais); zero fallbacks na versão final. Janela efetiva: 2025-04-01 (após 90 candles de \textit{warm-up}) até o último horizonte de cada par.

\begin{table*}[t]
  \centering
  \caption{Experimento 2 — Gemini Real 1d (252 inferências)}
  \label{tab:exp2_gemini}
  \small
  \begin{tabular}{lccccccc}
    \toprule
    Par & Amostras & Comp./Vend./Mant. & Acur. $t+1$ & Acur. prazo & Alvo & Ret. Gemini & B\&H (mesma janela) \\
    \midrule
    BTC/USDT & 63 & 38 / 7 / 18 & 48.9\% & \textbf{62.2\%} & 39.5\% & -30.97\% & -31.16\% \\
    ETH/USDT & 63 & 38 / 10 / 15 & 54.2\% & \textbf{54.2\%} & 47.4\% & -54.95\% & -17.48\% \\
    SOL/USDT & 63 & 39 / 9 / 15 & 54.2\% & \textbf{60.4\%} & 48.7\% & -64.07\% & -46.23\% \\
    DOGE/USDT& 63 & 43 / 6 / 14 & 51.0\% & 49.0\% & 32.6\% & -57.46\% & -56.25\% \\
    \bottomrule
  \end{tabular}
\end{table*}

Critérios do TCC no Experimento 2:

\begin{itemize}
  \item Acurácia de prazo > 50\%: \textbf{Atendido} — 3/4 ativos (BTC 62.2\%, ETH 54.2\%, SOL 60.4\%; DOGE 49.0\%).
  \item Retorno não negativo: \textbf{Não atendido} — 0/4 (todos negativos).
  \item \textbf{Conjunto dos dois critérios: Não atendido.}
\end{itemize}

Leitura: o Gemini real teve acurácia direcional de prazo acima de 50\% em 3 dos 4 ativos, mas essa acurácia não se converteu em retorno positivo no período de forte queda. Contra Buy-and-Hold na mesma janela, superou apenas BTC (+0.19 p.p.); ficou abaixo em ETH (-37.47 p.p.), SOL (-17.84 p.p.) e DOGE (-1.21 p.p.). A concordância com o proxy ficou entre 25.4\% e 42.9\%, mostrando que regras fixas de MM/RSI/MACD não reproduzem as decisões multimodais. O retorno é composto por decisão amostral; com passo de 7 dias e horizonte de 7–14 dias, há sobreposição. Não é curva de capital sem sobreposição e não inclui custos.

\subsection{Experimento 3 — Complemento Exploratório Intraday (1h/4h)}

Modelo \texttt{gemini-3.5-flash-lite}; 304 inferências (160 em 1h, 144 em 4h); zero fallbacks. Período: 01/07/2026 a 31/10/2026, amostragem a cada 3 dias, 90 candles de contexto. Primeira decisão após \textit{warm-up}: =04/07 em 1h, =16/07 em 4h.

\begin{table*}[t]
  \centering
  \caption{Experimento 3 — Gemini 1h/4h (304 inferências, jul–out/2026)}
  \label{tab:exp3_intraday}
  \small
  \begin{tabular}{lcccccccc}
    \toprule
    TF & Par & Amostras & Acur. $t+1$ & Acur. prazo & Alvo & Ret. Gemini & B\&H (mesma janela) \\
    \midrule
    1h & BTC/USDT & 40 & 64.0\% & \textbf{64.0\%} & 5.3\%  & -0.71\% & +23.00\% \\
    1h & ETH/USDT & 40 & 56.0\% & 48.0\% & 10.5\% & -1.89\% & +24.31\% \\
    1h & SOL/USDT & 40 & 56.7\% & \textbf{70.0\%} & 15.8\% & \textbf{+4.55\%} & +17.77\% \\
    1h & DOGE/USDT& 40 & 66.7\% & \textbf{61.1\%} & 0.0\%  & -0.81\% & \textbf{-4.62\%} \\
    4h & BTC/USDT & 36 & 45.5\% & \textbf{59.1\%} & 29.4\% & \textbf{+5.35\%} & +20.12\% \\
    4h & ETH/USDT & 36 & 61.9\% & \textbf{57.1\%} & 41.2\% & \textbf{+11.22\%} & +17.31\% \\
    4h & SOL/USDT & 36 & 37.0\% & \textbf{55.6\%} & 10.5\% & -2.08\%  & +24.73\% \\
    4h & DOGE/USDT& 36 & 33.3\% & \textbf{53.3\%} & 8.3\%  & -9.49\%  & +2.30\%  \\
    \bottomrule
  \end{tabular}
\end{table*}

\textbf{Importante:} o ``7 de 8'' significa 7 das 8 combinações par × timeframe; \textbf{não} significa 87,5\% de acertos. Agregando todos os pontos direcionais: acurácia de prazo = 108/183 = 59,0\% (1h: 60/98 = 61,2\%; 4h: 48/85 = 56,5\%) e acurácia $t+1$ = 97/183 = 53,0\%.

Critérios do TCC no Experimento 3:

\begin{table}[H]
  \centering
  \caption{Critérios TCC no Experimento 3}
  \label{tab:exp3_criterios}
  \small
  \begin{tabular}{lccc}
    \toprule
    Timeframe & Acurácia prazo > 50\% & Retorno não negativo & Conjunto \\
    \midrule
    1h & \textbf{Atendido} — 3/4 (BTC, SOL, DOGE) & Não — 1/4 (SOL +4.55\%) & Não \\
    4h & \textbf{Atendido} — 4/4 & Não — 2/4 (BTC +5.35\%, ETH +11.22\%) & Não \\
    \bottomrule
  \end{tabular}
\end{table}

Leitura: a acurácia de prazo superou 50\% em 7 de 8 combinações; o retorno não foi negativo em 3 de 8; Buy-and-Hold foi superado em apenas 1 de 8 (DOGE 1h: -0.81\% vs. -4.62\%). Em média simples, Gemini rendeu +0.77\% contra +15.62\% do Buy-and-Hold. O período foi de alta para BTC, ETH e SOL, favorecendo a estratégia passiva \textit{long-only}. A taxa de alvo ficou entre 0.0\% e 41.2\% (em 1h: 0.0\%–15.8\%), indicando projeções pouco realistas nesse timeframe. A janela de 4 meses é curta e não deve ser generalizada nem comparada diretamente ao Experimento 2 de 18 meses.

\subsection{Comparação de Estratégias — Janela Comum (Out-of-Sample LSTM)}

A Tabela~\ref{tab:comparacao_1d} mostra a janela comum do timeframe diário (2025-12-13 a 2026-06-29).

\begin{table*}[t]
  \centering
  \caption{Comparação na janela comum — 1d}
  \label{tab:comparacao_1d}
  \small
  \begin{tabular}{lccccc}
    \toprule
    Par & CriptoIA ret./acur. & Buy-and-Hold & LSTM ret./acur. & Baseline LSTM \\
    \midrule
    BTC/USDT & +1.28\% / 54.0\% & -35.03\% & -25.37\% / 52.3\% & 53.8\% \\
    ETH/USDT & +2.25\% / 46.5\% & -49.53\% & -48.59\% / 46.7\% & 53.8\% \\
    SOL/USDT & -14.05\% / 43.3\% & -44.67\% & -34.11\% / 47.2\% & 54.8\% \\
    DOGE/USDT& +15.26\% / 56.4\% & -48.16\% & -5.27\%  / 55.8\% & 59.3\% \\
    \bottomrule
  \end{tabular}
\end{table*}

Na grade completa 4×4 (16 combinações):
\begin{itemize}
  \item CriptoIA superou Buy-and-Hold em retorno em \textbf{15 de 16} combinações.
  \item LSTM teve acurácia > 50\% em \textbf{10 de 16}, mas superou o baseline de classe majoritária em apenas \textbf{6 de 16}; parte da acurácia é explicada por desbalanceamento de classes.
  \item CriptoIA teve acurácia > 50\% em \textbf{3 de 16} combinações.
  \item Não houve erro de execução do LSTM nas 16 combinações.
  \item Resultados semanais têm amostra muito pequena (28 válidos) e não devem ser generalizados.
  \item Não houve teste de significância estatística; números são descritivos.
\end{itemize}

\subsection{Testes e Cobertura}

\begin{itemize}
  \item \textbf{200 testes passed} em 21,76 s.
  \item \textbf{Cobertura total: 92\%} (backend + frontend + webapp; 1.660 statements, 126 não cobertos).
  \item Módulos com 100\%: \texttt{analise\_tecnica.py}, \texttt{gerador\_grafico.py}, \texttt{multiplos\_ativos.py}, \texttt{gerador\_dashboard.py}.
  \item Cobertura por módulo detalhada no RELATORIO\_TECNICO.md.
  \item CT01–CT06 todos cobertos e passando; testes de integração mockam rede e usam \texttt{tmp\_path}.
\end{itemize}

\subsection{Interface Web — Smoke Test}

Executado com cliente de teste Flask e CSRF ativo:

\begin{verbatim}
GET  /login       200
POST /login       302
GET  /            200
GET  /analise     200
GET  /historico   200
GET  /backtesting 200
POST /analise (sem CSRF) 400
GET  /logout      302
GET  /backtesting (após logout) 302
\end{verbatim}

Todas as rotas protegidas redirecionam corretamente; login válido/inválido, CSRF, headers e rate limit cobertos por testes.

\section{Conclusão}
\label{sec:conclusao}

Este trabalho entregou a etapa final do CriptoIA com:

\begin{enumerate}
  \item \textbf{Três experimentos de backtesting complementares e reproduzíveis:}
    \begin{itemize}
      \item Experimento 1: proxy determinístico 4×4 (16 combinações, 18 meses) — acurácia > 50\% em 2/4 ativos (1d), retorno não negativo em 1/4.
      \item Experimento 2: Gemini multimodal real em 1d amostral (252 inferências, zero fallbacks) — acurácia de prazo > 50\% em 3/4, retorno negativo em 4/4.
      \item Experimento 3: complemento exploratório Gemini em 1h/4h (304 inferências, zero fallbacks, 4 meses) — acurácia > 50\% em 7/8 combinações, retorno não negativo em 3/8, superação do Buy-and-Hold em 1/8.
    \end{itemize}
  \item \textbf{Comparação sistemática} com Buy-and-Hold e LSTM real (PyTorch, walk-forward, janela comum) — CriptoIA superou B\&H em 15/16 combinações; LSTM superou baseline majoritária em apenas 6/16.
  \item \textbf{Interface web Flask segura} com autenticação, CSRF, rate limit, histórico atômico e visualização dinâmica dos 3 experimentos.
  \item \textbf{Suite de 200 testes} com 92\% de cobertura, mocks de rede, sem dependência externa.
  \item \textbf{Metodologia transparente} com limitações explícitas (sobreposição de horizontes, long-only, sem custos, amostra semanal pequena, modelo gratuito, sem testes de significância).
\end{enumerate}

O \textbf{conjunto dos critérios de sucesso do TCC não foi atingido em nenhum dos três recortes}. O Gemini multimodal real demonstrou acurácia direcional de prazo superior a 50\% na maioria dos ativos (3/4 no 1d oficial; 7/8 no complemento 1h/4h), mas essa acurácia não se converteu consistentemente em retorno positivo nem em superação do Buy-and-Hold passivo nas mesmas janelas. A concordância entre o proxy determinístico e as decisões do Gemini ficou consistentemente baixa (25–65\%), confirmando que a IA não replica regras fixas de indicadores — porém, no período avaliado, a vantagem interpretativa não superou a estratégia passiva em mercado de queda (Exp. 2) nem em mercado de alta de curto prazo (Exp. 3).

A contribuição acadêmica central está em \textbf{reportar resultados negativos de forma transparente, com metodologia reproduzível e limitações declaradas}, em vez de apresentar números fabricados ou selecionados. O código-fonte completo, dados, relatórios e instruções de reprodução estão disponíveis em: \url{https://github.com/Teixeirx/tcc2}.

\section*{Agradecimentos}

O autor agradece ao Prof. Mirkos O. Martins, orientador deste trabalho, pelo suporte, direção e incentivo ao longo de ambas as etapas do TFG. Agradece também ao corpo docente do curso de Ciência da Computação da Universidade Franciscana e aos colegas que contribuíram com discussões técnicas durante o desenvolvimento.

\begin{thebibliography}{10}

\bibitem{nakamoto2008bitcoin}
S.~Nakamoto,
``Bitcoin: A Peer-to-Peer Electronic Cash System,''
2008. [Online]. Available: \url{https://bitcoin.org/bitcoin.pdf}

\bibitem{coinmarketcap2024}
CoinMarketCap,
``Global Cryptocurrency Market Cap Charts,''
2024. [Online]. Available: \url{https://coinmarketcap.com/charts/}

\bibitem{pinheiro2019mercado}
J.~L. Pinheiro,
\textit{Mercado de Capitais: Fundamentos e Técnicas}, 8.~ed.
São Paulo: Atlas, 2019.

\bibitem{assaf2012mercado}
A.~Assaf~Neto,
\textit{Mercado Financeiro}, 11.~ed.
São Paulo: Atlas, 2012.

\bibitem{murphy1999technical}
J.~J. Murphy,
\textit{Technical Analysis of the Financial Markets}.
New York: New York Institute of Finance, 1999.

\bibitem{cavalcante2005mercado}
F.~Cavalcante, F.~Misumi e R.~Rudge,
\textit{Mercado de Capitais: O Que É, Como Funciona}, 7.~ed.
Rio de Janeiro: Elsevier, 2005.

\bibitem{google2024gemini}
Google DeepMind,
``Gemini: A Family of Highly Capable Multimodal Models,''
\textit{arXiv preprint arXiv:2312.11805}, 2024.

\bibitem{maia2024inteligencia}
R.~Maia e P.~Gonçalves,
``Inteligência Artificial Generativa Aplicada a Finanças: Uma Revisão Sistemática,''
\textit{Revista Brasileira de Finanças}, vol.~22, no.~1, pp.~1--24, 2024.

\bibitem{edwards2018technical}
R.~D. Edwards, J.~Magee e W.~H.~C. Bassetti,
\textit{Technical Analysis of Stock Trends}, 11.~ed.
Boca Raton: CRC Press, 2018.

\bibitem{wilder1978new}
J.~W. Wilder,
\textit{New Concepts in Technical Trading Systems}.
Greensboro: Trend Research, 1978.

\bibitem{bollinger2001bollinger}
J.~Bollinger,
\textit{Bollinger on Bollinger Bands}.
New York: McGraw-Hill, 2001.

\bibitem{liu2023visual}
H.~Liu, C.~Li, Q.~Wu e Y.~J. Lee,
``Visual Instruction Tuning,''
in \textit{Proc. NeurIPS}, 2023.

\bibitem{binance2024}
Binance,
``Binance API Documentation,''
2024. [Online]. Available: \url{https://binance-docs.github.io/apidocs/}

\bibitem{ccxt2024}
CCXT,
``CCXT — CryptoCurrency eXchange Trading Library,''
2024. [Online]. Available: \url{https://github.com/ccxt/ccxt}

\bibitem{fischer2018deep}
T.~Fischer e C.~Krauss,
``Deep learning with long short-term memory networks for financial market predictions,''
\textit{European Journal of Operational Research}, vol.~270, no.~2, pp.~654--669, 2018.

\bibitem{cavalcante2015computational}
R.~C. Cavalcante \textit{et al.},
``Computational Intelligence and Financial Markets: A Survey and Future Directions,''
\textit{Expert Systems with Applications}, vol.~55, pp.~194--211, 2016.

\bibitem{monteiro2026ssinnova}
R.~P.~N. Santos e C.~S. Capriglione,
``Sistema de Análise de Criptomoedas em Tempo Real para Apoio à Decisão de Investidores,''
\textit{RECIMA21 --- Revista Científica Multidisciplinar}, vol.~7, n.~5, 2026.
DOI: \url{https://doi.org/10.47820/recima21.v7i5.7927}

\bibitem{silvajunior2022allbot}
D.~S. Silva~Júnior, G.~P. Silva e R.~N.~O. Muniz,
``AllBot: plataforma de robôs para investimento no mercado de criptomoedas,''
Trabalho de Conclusão de Curso (Bacharelado em Engenharia de Computação),
Instituto Federal de Educação, Ciência e Tecnologia da Paraíba (IFPB),
Campina Grande, 2022.
Available: \url{http://repositorio.ifpb.edu.br/handle/177683/2465}

\bibitem{franciscon2023lstm}
K.~F. Franciscon,
``Predição do preço da bitcoin por meio da incorporação de indicadores técnicos em modelos LSTM,''
Trabalho de Conclusão de Curso (Bacharelado em Engenharia de Computação),
Universidade Tecnológica Federal do Paraná (UTFPR), Toledo, 2023.
Available: \url{http://repositorio.utfpr.edu.br/jspui/handle/1/36696}

\bibitem{hochreiter1997lstm}
S.~Hochreiter e J.~Schmidhuber,
``Long Short-Term Memory,''
\textit{Neural Computation}, vol.~9, no.~8, pp.~1735--1780, 1997.

\bibitem{bailey2014pseudomathematics}
D.~H. Bailey, J.~M. Borwein, M.~L. de Prado e Q.~J. Zhu,
``Pseudo-Mathematics and Financial Charms,''
\textit{Notices of the AMS}, vol.~61, no.~5, pp.~458--471, 2014.

\bibitem{zhang2024finben}
Y.~Zhang \textit{et al.},
``FinBen: A Comprehensive Benchmark for Evaluating Large Language Models in Financial NLP,''
\textit{arXiv preprint arXiv:2402.12345}, 2024.

\bibitem{yang2024fingpt}
H.~Yang \textit{et al.},
``FinGPT: Open-Source Financial Large Language Models,''
\textit{arXiv preprint arXiv:2306.06031}, 2023.

\bibitem{lopez2024chatgpt}
A.~Lopez-Lira e Y.~Tang,
``Can ChatGPT Forecast Stock Price Movements? Return Predictability and Large Language Models,''
\textit{arXiv preprint arXiv:2304.07619}, 2023.

\bibitem{owasp2021csrf}
OWASP,
``Cross-Site Request Forgery (CSRF) Prevention Cheat Sheet,''
2021. [Online]. Available: \url{https://cheatsheetseries.owasp.org/cheatsheets/Cross-Site_Request_Forgery_Prevention_Cheat_Sheet.html}

\bibitem{owasp2021headers}
OWASP,
``Security Headers,''
2021. [Online]. Available: \url{https://owasp.org/www-project-secure-headers/}

\bibitem{tcc1}
G.~Teixeira,
``CriptoIA: Um Sistema de Análise de Criptomoedas Baseado em Indicadores Técnicos e Inteligência Artificial Multimodal,''
Trabalho de Conclusão de Curso I, Universidade Franciscana, 2026.

\end{thebibliography}

\end{document}